\textbf{COLUNI 2025 -- Questão 03}\par\medskip Leia o trecho abaixo.

\par

Seja para os trabalhadores da Fazenda Água Negra, seja para as pessoas em sofrimento em qualquer parte do Brasil ou do mundo, as vítimas de sistemas de opressão e exploração tendem a se identificar com os valores transmitidos por Torto Arado, sentindo simpatia pelos eventos enfrentados pelos protagonistas. É desta maneira que Torto Arado realizou a transição da “escrita de nós” para a “escrita da gente”, oferecendo-nos um novo modelo de narrativa coletiva. (XING, Fan. De “nós” para “a gente: a narrativa coletiva em                                Torto   arado.    Disponível     em: https://www.scielo.br/j/tla/a/gYXjKjj9SNkZNT54z5TNFNJ/. Acesso em: 08 jul. 2024.)

\par

Considerando o trecho acima e a obra \textit{Torto arado}, marque a afirmativa INCORRETA em relação às narradoras da obra.\begin{enumerate}
\item \textit{Torto arado} consegue, ao focar em personagens femininas, abraçar uma inclusividade que
       potencialmente une gêneros, raças, classes e regiões. Argumenta-se que a transição da “escrita de nós”
       para a “escrita da gente” simboliza um novo paradigma na narrativa coletiva.
\item À medida que Bibiana, a narradora da primeira parte de \textit{Torto arado}, “Fio de corte”, cresce, “nós” diversas
       vezes aponta aos arrendatários da fazenda, como em “falou sobre a escola que não seria suficiente para
       completarmos os estudos, mas que era um grande benefício para nós que morávamos em Água Negra,
       carentes de tudo” (VIEIRA JUNIOR, 2019, p. 71).
\item Na segunda parte do romance \textit{Torto arado}, “nós” tem múltiplas referências, como crença: “como seria
       tudo sem as forças dos encantados, que por tanto tempo haviam estado entre nós” (VIEIRA JUNIOR,
       2019, p. 194), apontando a importância da crença jarê para este grupo.
\item Em “Rio de sangue”, a terceira parte do romance, a coletividade/inclusividade é menos evidente que nas
       duas primeiras partes, visto que a narradora da terceira parte é Santa Rita Pescadeira, uma entidade do
       jarê, que já se abrigou em muitos corpos, não possuindo uma natureza coletiva.\end{enumerate}
